\newcommand{\plotPGFfile}{true} % PGF olarak DATA dosyaları çizilsin mi çizilmesin mi? Çizilmez ise dosyadan sadece data okunan kısım atlanıyor, resmin çerçevesi oluşturuluyor...
\definecolor{dkgreen}{rgb}{0,0.6,0}
\definecolor{gray}{rgb}{0.5,0.5,0.5}
\definecolor{mauve}{rgb}{0.58,0,0.82}

\lstset{frame=tb,
	language=Java, %Uygulamanın yazıldığı programlama dili, (Örnek kullanım chapter4.tex'te verilmiştir.)
	aboveskip=3mm,
	belowskip=3mm,
	showstringspaces=false,
	columns=flexible,
	basicstyle={\small\ttfamily},
	numbers=none,
	numberstyle=\tiny\color{gray},
	keywordstyle=\color{blue},
	commentstyle=\color{dkgreen},
	stringstyle=\color{mauve},
	breaklines=true,
	breakatwhitespace=true,
	tabsize=3
}

\author{Mustafa}{G{\"O}KTEPE}{1030510749}{1030510749@erciyes.edu.tr}{05331556738}%1. Öğrenci ad soyad
\authorIki{Ahmet}{ÇELİK}{1030510691}{1030510691@erciyes.edu.tr}{05453246440}%2. Öğrenci ad soyad Opsiyonel
\authorUc{Mustafa}{KÜPELİ}{1030510266}{1030510266@erciyes.edu.tr}{05360100620}%3. Öğrenci ad soyad  Opsiyonel
%\authorDort{Ad}{Soyad}{103051004}{numara4@erciyes.edu.tr}{054xxxxxxx}%4. Öğrenci ad soyad Opsiyonel

\anabilimdali{Bilgisayar Mühendisliği}
\anabilimdaliUpper{BİLGİSAYAR MÜHENDİSLİĞİ}

\date{Kasım 2026}
\Yil{2026-2027} %Eğitim Öğretim Yili
\Donem{Güz} %Donemi
\teziyoneten{Ar.\ G{\"O}r. HAYRİ ÖZTÜRK} % Tez Danismani





\title{MEDIAPIPE VE MAKİNE ÖĞRENMESİ KULLANARAK TÜRK İŞARET DİLİ (TİD) TABANLI GERÇEK ZAMANLI EL HAREKETİ TANIMA VE ÇEVİRİ SİSTEMİ} %TAMAMEN BÜYÜK HARF\vspace*{5mm}
\titleCase{MediaPipe ve Makine Öğrenmesi Kullanarak Türk İşaret Dili (TİD) Tabanlı Gerçek Zamanlı El Hareketi Tanıma ve Çeviri Sistemi} %Sadece ilk Harfleri Büyük


